\documentclass[10pt,a4paper]{amsart}
\usepackage[margin=19mm]{geometry}
\usepackage{amsmath,amssymb,mathtools}
\usepackage[scr=rsfs,cal=euler]{mathalfa}
\usepackage{xfrac}
\usepackage[unicode,hidelinks,bookmarks=false]{hyperref}
\newcommand{\ie}{i.e.\ }
\newcommand{\cf}{cf.\ }
\newcommand{\resp}{respectively\ }
\newcommand{\Subsection}{\subsection}
\providecommand{\nobreakdash}{\nobreak\hbox{-}}
\setlength{\emergencystretch}{2em}
\multlinegap0pt
\usepackage{amssymb}
\usepackage{bm}
\usepackage{enumerate}
\usepackage{xcolor}
\newtheorem{thm}{Theorem}
\newtheorem{pf}{Pf}
\newcommand{\N}{\mathbb N}
\newcommand{\R}{\mathbb R}
\newcommand{\dl}{\delta}
\newcommand{\ep}{\varepsilon}
\title{\bf From Tsallis to KL: Convergence and Error Estimates for Tsallis-Regularized Optimal Transport}
\author{Takeshi Suguro, Toshiaki Yachimura}
\date{Source: arXiv:2609.06432v1 (CC BY 4.0)}
\begin{document}
\maketitle

\noindent\emph{Source and licence.} \bf From Tsallis to KL: Convergence and Error Estimates for Tsallis-Regularized Optimal Transport. Version arXiv:2609.06432v1 (CC BY 4.0), distributed by arXiv under the Creative Commons Attribution 4.0 International licence, http://creativecommons.org/licenses/by/4.0/, as recorded in the arXiv licence metadata for this exact version and shown on the abstract page https://arxiv.org/abs/2609.06432v1. Full licence notice: \texttt{licenses/arxiv-2609.06432-CC-BY-4.0.txt}.\par\smallskip

\noindent\textit{Editorial scope.} Counted unit: the article's thm thm;narrow (source0.tex lines 394--396). The article's complete original proof of it is delivered with it (source0.tex lines 628--843, one \texttt{proof} environment of 216 lines, which the article places away from the statement). No mathematics is written, repaired or re-derived: the statement and the proof are the article's own lines, in the article's own notation, and no retained line is substituted or re-flowed.  No further source-local context is needed: every cross-reference of every retained line resolves inside the two delivered spans.  Nothing was omitted from any retained span, so the package carries no omission record.  No retained line contains a citation, so the wrapper carries no bibliography and no bibliography entry.  No retained line contains a figure environment, an image-inclusion command or drawing code, so no image is delivered and the package declares no asset dependency.  Editorial formatting only, in addition: an AMS \texttt{amsart} preamble that restates the article's own declarations and macros used by the retained lines, namely the environments \texttt{thm}, \texttt{pf} on separate counters, together with the standard packages those lines need.  The retained environments print in this document as Theorem 1, Pf 1 in order of appearance, whereas the article prints the same items with the numbering of its own full document.  The complete native source of the article as distributed in the arXiv e-print for this exact version is bundled unchanged as \texttt{sources/209535/source0.tex}.
\begin{thm}\label{thm;narrow}
Let $X_i$ be Polish spaces, let $\mu_i\in\mathcal P(X_i)$ $(i=1,2)$, and let $c\in C_b(X_1\times X_2)$ be nonnegative. Let $\ep>0$ and $\{q_k\}_{k\in\N}\subset(1,2]$ satisfy $q_k\downarrow1$. Then $\{F_k\}_{k\in\N}$ $\Gamma$-converges to $F$ with respect to the narrow topology. Moreover, the unique minimizer $\pi_k^*$ of $F_k$ converges narrowly to the unique minimizer $\pi^*$ of $F$.
\end{thm}

\begin{pf}{Theorem \ref{thm;narrow}}
Let $\{q_k\}_{k \in \N} \subset (1,2]$ satisfy $q_k\downarrow1$. We define the functionals by
\begin{equation*}
F_k(\pi) = \left\{
\begin{aligned} &\int_{X_1 \times X_2} c\, d\pi + \ep D_{q_k}(\pi, m)
&& \mathrm{for}\ \pi \in \Pi(\mu_1, \mu_2), \\
&+ \infty, && \mathrm{otherwise},
\end{aligned}
\right.
\end{equation*}
and
\begin{equation*}
F(\pi) = \left\{
\begin{aligned}
&\int_{X_1 \times X_2} c\, d\pi + \ep \mathrm{KL}(\pi, m)
&& \mathrm{for}\ \pi \in \Pi(\mu_1, \mu_2), \\
&+ \infty, && \mathrm{otherwise}.
\end{aligned}
\right.
\end{equation*}

First, we prove the liminf inequality
\begin{equation}\label{liminf;0}
  F(\pi)
  \leq \liminf_{k \to \infty} F_k(\pi_k)
\end{equation}
for any $\{\pi_k\}_{k \in \N}$ such that $\pi_k \to \pi$ narrowly. If $\displaystyle \liminf_{k\to\infty}F_k(\pi_k)=+\infty$, there is nothing to prove. Otherwise, after passing to a subsequence realizing the liminf, we may assume that $F_k(\pi_k)<+\infty$ for every $k$. Hence $\pi_k\in\Pi(\mu_1,\mu_2)$, and the narrow closedness of $\Pi(\mu_1,\mu_2)$ implies that $\pi\in\Pi(\mu_1,\mu_2)$. Moreover, for every probability measure $\eta$,
\begin{equation*}
\mathrm{KL}(\eta,m) \leq D_{q_k}(\eta,m).
\end{equation*}
Indeed, when $\eta \ll m$, this follows from $\log x\leq\log_{2-q_k}x$ for $x>0$, while otherwise both sides are $+\infty$. Consequently, the lower semicontinuity of the Kullback--Leibler divergence gives
\begin{equation}\label{liminf;1}
  \ep \mathrm{KL}(\pi, m)
  \leq \liminf_{k \to \infty} \ep \mathrm{KL}(\pi_k, m)
  \leq \liminf_{k \to \infty} \ep D_{q_k}(\pi_k, m).
\end{equation}
On the other hand, by the lower semicontinuity of the transport cost functional,
\begin{equation}\label{liminf;2}
\int_{X_1 \times X_2} c\, d\pi \leq \liminf_{k \to \infty} \int_{X_1 \times X_2} c\, d\pi_k.
\end{equation}
Combining \eqref{liminf;1} and \eqref{liminf;2}, and using the superadditivity of the $\liminf$, we obtain
\begin{equation*}
F(\pi) \leq \liminf_{k\to\infty} F_k(\pi_k),
\end{equation*}
which proves \eqref{liminf;0}.

Next, we prove the limsup inequality. For any $\pi\in\mathcal P(X_1\times X_2)$, we seek a sequence $\{\pi_k\}_{k\in\N}$ converging narrowly to $\pi$ such that
\begin{equation}\label{limsup;0}
  F(\pi) \geq \limsup_{k \to \infty} F_k(\pi_k).
\end{equation}
It is enough to consider the case $F(\pi)<+\infty$, since otherwise the constant sequence $\pi_k = \pi$ satisfies the limsup inequality. Thus $\pi\in\Pi(\mu_1,\mu_2)$ and $d\pi=\rho\,dm$ for some density $\rho$.
For $n \in \N$, we set $\tilde{\rho}_n = \min\{\rho, n\}$, 
\begin{equation*}
  \dl_n = 1 - \int_{X_1 \times X_2} \tilde{\rho}_n\, dm,
\end{equation*}
\begin{equation*}
a_n(x) = 1 - \int_{X_2} \tilde{\rho}_n(x, y)\, d\mu_2(y), \quad \mathrm{and} \quad b_n(y) = 1 - \int_{X_1} \tilde{\rho}_n(x, y)\, d\mu_1(x). 
\end{equation*}
We note that
\begin{equation*}
\int_{X_1} a_n(x)\, d\mu_1 = \int_{X_2} b_n(y)\, d\mu_2 = \dl_n.
\end{equation*}
We set 
\begin{equation*}
\rho_n(x, y) = \left\{
    \begin{aligned}
    &\tilde{\rho}_n(x, y) + \frac{a_n(x) b_n(y)}{\dl_n}
    && \mathrm{if}\ \dl_n > 0, \\
    &\rho
    && \mathrm{if}\ \dl_n = 0
    \end{aligned}
    \right.
\end{equation*}
and $d\pi_n =\rho_n\,dm$. Then $\pi_n \in \Pi(\mu_1, \mu_2)$. If $\dl_n>0$, since $\dl_n\to0$ by monotone convergence,
\begin{equation*}
\|\pi_n-\pi\|_{\mathrm{TV}} = \int_{X_1\times X_2}|\rho_n-\rho|\,dm \leq \int_{X_1\times X_2} |\rho-\tilde\rho_n| \,dm +\int_{X_1\times X_2}\frac{a_nb_n}{\dl_n}\,dm = 2\dl_n\to0. 
\end{equation*}
If $\dl_n=0$, then $\rho_n=\rho$ by definition. Therefore, $\pi_n\to\pi$ in total variation, and hence narrowly.
Moreover, since $\int_{X_1\times X_2}|\rho_n-\rho|\,dm\to0$, we have
\begin{equation}\label{limsup;2}
\left|\int_{X_1\times X_2}c\,d\pi_n-\int_{X_1\times X_2}c\,d\pi\right|
\leq \|c\|_\infty\int_{X_1\times X_2}|\rho_n-\rho|\,dm\to0.
\end{equation}
Furthermore, by the marginal conditions for $\rho$ and the inequality $\tilde{\rho}_n\leq\rho$, we have $0\leq a_n(x)\leq1$ for $\mu_1$-a.e. $x\in X_1$ and $0\leq b_n(y)\leq1$ for $\mu_2$-a.e. $y\in X_2$. Therefore, if $\dl_n>0$, then
\begin{equation}\label{limsup;1}
\rho_n(x,y)\leq\tilde{\rho}_n(x,y)+\frac{1}{\dl_n}\leq n+\frac{1}{\dl_n}\quad m\text{-a.e.}
\end{equation}
Since $\dl_n>0$, the right-hand side is finite for each fixed $n$, and hence $\rho_n\in L^\infty(m)$. If $\dl_n=0$, then
\begin{equation*}
0=\dl_n=\int_{X_1\times X_2}(\rho-\tilde{\rho}_n)\,dm.
\end{equation*}
Since $\rho-\tilde{\rho}_n\geq0$, it follows that $\rho=\tilde{\rho}_n$ $m$-a.e. By the definition of $\rho_n$, we then have
\begin{equation*}
\rho_n=\rho=\tilde{\rho}_n\leq n\quad m\text{-a.e.}
\end{equation*}
Thus, in either case, $\rho_n\in L^\infty(m)$ for every $n$.

In what follows, we show 
\begin{equation}\label{limsup;3}
\lim_{n\to\infty}\mathrm{KL}(\pi_n,m)=\mathrm{KL}(\pi,m).
\end{equation}
If $\dl_n=0$ for some $n$, then $\rho=\tilde{\rho}_n\leq n$ $m$-a.e. Hence, for every $j\geq n$, we have $\tilde{\rho}_j=\rho$, and therefore $\dl_j=0$ and $\pi_j=\pi$. In this case, \eqref{limsup;3} follows immediately. Otherwise, $\dl_n>0$ for every $n$. We then set
\begin{equation*}
r_n = \frac{\tilde{\rho}_n}{1-\dl_n} \quad\mathrm{and}\quad s_n = \frac{a_nb_n}{\dl_n^2}.
\end{equation*}
Since
\begin{equation*}
\int_{X_1\times X_2}r_n\,dm = 1 \quad\mathrm{and}\quad \int_{X_1\times X_2}s_n\,dm = 1,
\end{equation*}
both $r_n$ and $s_n$ are probability densities with respect to $m$. Moreover,
\begin{equation*}
\rho_n = (1-\dl_n)r_n + \dl_n s_n.
\end{equation*}
Applying Jensen's inequality to the convex function $t\mapsto t\log t$, we obtain
\begin{align*}
\mathrm{KL}(\pi_n,m)
&=\int_{X_1\times X_2}\rho_n\log\rho_n\,dm \\
&\leq (1-\dl_n)\int_{X_1\times X_2}r_n\log r_n\,dm
+\dl_n\int_{X_1\times X_2}s_n\log s_n\,dm \\
&=\int_{X_1\times X_2}\tilde{\rho}_n\log\tilde{\rho}_n\,dm
-(1-\dl_n)\log(1-\dl_n) \\
&\quad +\int_{X_1}a_n\log a_n\,d\mu_1
+\int_{X_2}b_n\log b_n\,d\mu_2
-2\dl_n\log\dl_n.
\end{align*}
Since $\tilde{\rho}_n\leq\rho$, we have
\begin{equation*}
(\tilde{\rho}_n\log\tilde{\rho}_n)_+ \leq (\rho\log\rho)_+.
\end{equation*}
Moreover, $t\log t\geq-e^{-1}$ for every $t\geq0$, where $0\log0=0$. Hence
\begin{equation*}
\left|\tilde{\rho}_n\log\tilde{\rho}_n\right| \leq (\rho\log\rho)_+ + e^{-1} \quad m\text{-a.e.}
\end{equation*}
Since $\mathrm{KL}(\pi,m)<+\infty$, the right-hand side is integrable with respect to $m$. Therefore, by the dominated convergence theorem, we have
\begin{equation*}
\lim_{n\to\infty}\int_{X_1\times X_2}\tilde{\rho}_n\log\tilde{\rho}_n\,dm = \int_{X_1\times X_2}\rho\log\rho\,dm.
\end{equation*}
Furthermore, by the marginal identities for $\rho$ and the monotone convergence theorem,
\begin{equation*}
\int_{X_2}\tilde{\rho}_n(x,y)\,d\mu_2(y)\uparrow1
\quad\text{for $\mu_1$-a.e. }x,
\end{equation*}
and
\begin{equation*}
\int_{X_1}\tilde{\rho}_n(x,y)\,d\mu_1(x)\uparrow1
\quad\text{for $\mu_2$-a.e. }y.
\end{equation*}
Hence $a_n\downarrow0$ $\mu_1$-a.e. and $b_n\downarrow0$ $\mu_2$-a.e. Since $0\leq a_n,b_n\leq1$ and $|t\log t|\leq e^{-1}$ on $[0,1]$, where $0\log0=0$, the dominated convergence theorem gives
\begin{equation*}
\lim_{n\to\infty}\int_{X_1}a_n\log a_n\,d\mu_1=0 \quad\mathrm{and}\quad
\lim_{n\to\infty}\int_{X_2}b_n\log b_n\,d\mu_2=0.
\end{equation*}
Since $\dl_n\to0$, we also have
\begin{equation*}
(1-\dl_n)\log(1-\dl_n)\to0 \quad\mathrm{and}\quad \dl_n\log\dl_n\to0.
\end{equation*}
Therefore,
\begin{equation*}
\limsup_{n\to\infty}\mathrm{KL}(\pi_n,m)\leq\mathrm{KL}(\pi,m).
\end{equation*}
On the other hand, since $\pi_n\to\pi$ narrowly and the Kullback--Leibler divergence is lower semicontinuous with respect to narrow convergence,
\begin{equation*}
\liminf_{n\to\infty}\mathrm{KL}(\pi_n,m)\geq\mathrm{KL}(\pi,m).
\end{equation*}
Consequently, we obtain \eqref{limsup;3}.

For every fixed $n$, the density $\rho_n$ is bounded. Let $M_n<+\infty$ be such that $0\leq\rho_n\leq M_n$ $m$-a.e. For $t>0$ and $1<q\leq2$, the mean value theorem applied to the function $s\mapsto t^s$ yields
\begin{equation*}
\frac{t^q-t}{q-1}=t^\xi\log t
\end{equation*}
for some $\xi\in(1,q)$. At points where $\rho_n=0$, both $(\rho_n^q-\rho_n)/(q-1)$ and $\rho_n\log\rho_n$ are equal to zero. Hence
\begin{equation*}
\frac{\rho_n^q-\rho_n}{q-1}\to\rho_n\log\rho_n \quad m\text{-a.e. as }q \downarrow 1.
\end{equation*}
Moreover, if $0<t\leq1$, then $|t^\xi\log t|\leq t|\log t|\leq e^{-1}$, while if $1\leq t\leq M_n$, then $|t^\xi\log t|\leq M_n^2\log(\max\{M_n,1\})$. Thus the integrands are dominated by an $m$-integrable constant independent of $q$. Therefore, by the dominated convergence theorem,
\begin{equation*}
D_q(\pi_n,m)\to\mathrm{KL}(\pi_n,m)
\quad\text{as }q\downarrow1.
\end{equation*}

We now choose a diagonal sequence so that $n\to\infty$ and $q_k\downarrow1$ simultaneously. For each $n\in\N$, there exists $K_n\in\N$ such that
\begin{equation*}
|D_{q_k}(\pi_n,m)-\mathrm{KL}(\pi_n,m)|<\frac{1}{n}
\end{equation*}
for every $k\geq K_n$. Increasing $K_n$ if necessary, we may assume that $K_1<K_2<\cdots$. Define
\begin{equation*}
n(k)=
\begin{cases}
1, & k<K_1,\\
\max\{n\in\N:K_n\leq k\}, & k\geq K_1.
\end{cases}
\end{equation*}
Then $n(k)\to\infty$ as $k\to\infty$. Set $\hat{\pi}_k=\pi_{n(k)}$. Since $n(k)\to\infty$ and $\pi_n\to\pi$ narrowly, we have $\hat{\pi}_k\to\pi$ narrowly. Moreover, for every $k\geq K_1$, we have $k\geq K_{n(k)}$, and hence
\begin{equation*}
|D_{q_k}(\hat{\pi}_k,m)-\mathrm{KL}(\hat{\pi}_k,m)|<\frac{1}{n(k)}.
\end{equation*}
Therefore,
\begin{equation*}
D_{q_k}(\hat{\pi}_k,m)\leq\mathrm{KL}(\hat{\pi}_k,m)+\frac{1}{n(k)}.
\end{equation*}
Since $n(k)\to\infty$, \eqref{limsup;3} gives
\begin{equation}\label{limsup;4}
\limsup_{k\to\infty}D_{q_k}(\hat{\pi}_k,m)\leq\mathrm{KL}(\pi,m).
\end{equation}
Finally, \eqref{limsup;2} also implies
\begin{equation*}
\int_{X_1\times X_2}c\,d\hat{\pi}_k\to\int_{X_1\times X_2}c\,d\pi.
\end{equation*}
Combining this convergence with \eqref{limsup;4} yields \eqref{limsup;0}.

Recall that the sequence $\{F_k\}$ is equicoercive on $\mathcal{P}(X_1\times X_2)$ if, for every $M\in\R$, there exists a narrowly compact set $K_M\subset\mathcal{P}(X_1\times X_2)$ such that
\begin{equation*}
\{\pi\in\mathcal{P}(X_1\times X_2):F_k(\pi)\leq M\}\subset K_M
\end{equation*}
for every $k\in\N$. Since $F_k(\pi)<+\infty$ implies $\pi\in\Pi(\mu_1,\mu_2)$, every sublevel set of $F_k$ is contained in $\Pi(\mu_1,\mu_2)$. Since $\Pi(\mu_1,\mu_2)$ is narrowly compact, we may take $K_M=\Pi(\mu_1,\mu_2)$ for every $M\in\R$. Hence the sequence $\{F_k\}$ is equicoercive. Each $F_k$, as well as $F$, is lower semicontinuous with respect to narrow convergence. Moreover, $m\in\Pi(\mu_1,\mu_2)$ has finite energy for all $F_k$ and $F$, and hence minimizers exist. Strict convexity of the entropy terms on the convex set $\Pi(\mu_1,\mu_2)$ gives uniqueness. By the fundamental theorem of $\Gamma$-convergence and the equicoercivity of $\{F_k\}$, every cluster point of $\{\pi_k^*\}_{k\in\N}$ is a minimizer of $F$. Since $F$ has the unique minimizer $\pi^*$, every cluster point coincides with $\pi^*$. Hence the whole sequence $\{\pi_k^*\}_{k\in\N}$ converges narrowly to $\pi^*$.
\end{pf}

\end{document}
